% Generated by make_fig_b_architecture.py. Do not edit by hand.
\documentclass[tikz,border=2pt]{standalone}
\usepackage[T1]{fontenc}
\usepackage{newtxtext,newtxmath}
\usepackage[varqu,scaled=0.95]{zi4}
\usetikzlibrary{arrows.meta,positioning,calc,fit,backgrounds}
\definecolor{ink}{HTML}{1F2328}\definecolor{mute}{HTML}{57606A}
\definecolor{newC}{HTML}{24527F}\definecolor{extC}{HTML}{D8DEE4}\definecolor{lane}{HTML}{F3F5F7}
\definecolor{exNew}{HTML}{DCE7F2}
\tikzset{
  >={Stealth[length=5.5pt,width=4.6pt]},
  box/.style={draw=ink, line width=0.6pt, rounded corners=2pt, align=center, inner sep=3pt,
              font=\fontsize{7.6}{9}\selectfont, text=ink, minimum height=1.75cm, anchor=center},
  reused/.style={box, fill=white},
  ext/.style={box, fill=extC},
  new/.style={box, fill=newC, text=white, draw=newC},
  io/.style={box, dashed, fill=white},
  arr/.style={->, line width=0.9pt, draw=ink},
  lab/.style={font=\fontsize{6.6}{7.5}\selectfont\itshape, text=mute, inner sep=1pt},
  lanet/.style={font=\fontsize{8}{9}\selectfont\bfseries, text=mute, anchor=north},
  key/.style={font=\fontsize{7.4}{9}\selectfont, text=ink, anchor=west},
}
% a box: title, what it does, and the running example on a strip
\newcommand{\bx}[3]{\textbf{#1}\\[1pt]{#2}\\[3pt]\setlength{\fboxsep}{2pt}\colorbox{white}{\parbox{\dimexpr\linewidth-4pt\relax}{\centering\color{ink}\fontsize{6.6}{8}\selectfont #3}}}
\newcommand{\bxn}[3]{\textbf{#1}\\[1pt]{#2}\\[3pt]\setlength{\fboxsep}{2pt}\colorbox{exNew}{\parbox{\dimexpr\linewidth-4pt\relax}{\centering\color{ink}\fontsize{6.6}{8}\selectfont #3}}}
\begin{document}
\begin{tikzpicture}
\node[io, text width=2.0cm] (qasm) at (1.0,0.0) {\bx{OpenQASM 2, 3}{\texttt{.qasm} file}{\texttt{cx q[0],q[1];}}};
\node[io, text width=2.0cm] (py) at (1.0,-2.3) {\bx{Qiskit}{\texttt{.py} program}{\texttt{qc.cx(q[src], q[a])}}};
\node[new, text width=2.85cm] (fe) at (4.085,0.0) {\bxn{Front ends}{OpenQASM~2 and~3 to srcML, with positions}{\texttt{\textless{}argument\_list type=\textquotedbl{}quantum\textquotedbl{}\textgreater{}}}};
\node[reused, text width=2.85cm] (srcml) at (4.085,-2.3) {\bx{srcML}{Python markup with positions}{\texttt{\textless{}call pos:start=\textquotedbl{}14:5\textquotedbl{}\textgreater{}}}};
\node[ext, text width=2.85cm] (ss) at (4.085,-4.6) {\bx{srcSlice, patched}{def, use, dvars, cfuncs of each variable}{\texttt{src}: def 9, 22}};
\node[new, text width=2.65cm] (e1) at (7.495,-2.3) {\bxn{E1 Roles}{what each call does to which qubits}{\texttt{cx(control, target)}}};
\node[new, text width=2.65cm] (e2) at (10.805,-2.3) {\bxn{E2 Trace}{the circuit, built without running}{line 14, $h=1$: \texttt{cx q2, q3}}};
\node[new, text width=2.65cm] (e3) at (14.115,-2.3) {\bxn{E3 Slice}{light cone through measured bits and loops}{22 of 34 operations}};
\node[new, text width=2.65cm] (e4) at (14.115,-4.6) {\bxn{E4 Classical lines}{srcSlice's closure from the slice's lines}{\texttt{src = b}, line 22}};
\node[new, text width=2.65cm] (e5) at (14.115,0.0) {\bxn{E5 Profile}{uses, names and dependents of a qubit}{\texttt{q[b]}, \texttt{q[src]}}};
\node[io, text width=2.8cm] (qm) at (17.499999999999996,2.05) {\bx{Quantum markup}{srcML with the roles and operations, for srcML tools}{\texttt{qu:role=\textquotedbl{}gate\textquotedbl{} qu:gate=\textquotedbl{}cx\textquotedbl{}}}};
\node[io, text width=2.8cm] (o1) at (17.499999999999996,0.0) {\bx{Qubit profiles}{\texttt{q[4]} depends on 6 qubits}{names \texttt{q[b]}, \texttt{q[src]}}};
\node[io, text width=2.8cm] (o2) at (17.499999999999996,-2.3) {\bx{Slice}{source lines and iterations}{lines 8, 12--17, 19, 21; $h=0$--2}};
\node[io, text width=2.8cm] (o3) at (17.499999999999996,-4.6) {\bx{Classical lines}{the statements that set the qubits}{lines 3--7, 9--11, 18, 20, 22}};
\begin{scope}[on background layer]
\fill[lane, rounded corners=3pt] (-0.200,3.070) rectangle (2.200,-5.620);
\node[lanet] at (1.000,3.490) {Input};
\fill[lane, rounded corners=3pt] (2.460,3.070) rectangle (5.710,-5.620);
\node[lanet] at (4.085,3.490) {Markup};
\fill[lane, rounded corners=3pt] (5.970,3.070) rectangle (15.640,-5.620);
\node[lanet] at (10.805,3.490) {Quantum layer};
\fill[lane, rounded corners=3pt] (15.900,3.070) rectangle (19.100,-5.620);
\node[lanet] at (17.500,3.490) {Output};
\end{scope}
\draw[arr] (qasm) -- (fe);
\draw[arr] (py) -- (srcml);
\draw[arr] (srcml) -- (e1);
\draw[arr] (fe.east) -| node[lab, pos=0.25, above] {markup} ([xshift=-0.55cm]e1.north);
\draw[arr] (e1) -- (e2);
\draw[arr] (e2) -- (e3);
\draw[arr] (e2.north) |- node[lab, pos=0.75, above] {roles and operations, as attributes} (qm.west);
\draw[arr] (e3) -- (e5);
\draw[arr] (e3) -- node[lab, right] {quantum lines} (e4);
\draw[arr] (srcml) -- (ss);
\draw[arr] (ss) -- node[lab, above] {profiles} (e4);
\draw[arr] (e5) -- (o1);
\draw[arr] (e3) -- (o2);
\draw[arr] (e4) -- (o3);
\node[reused, minimum height=0.3cm, minimum width=0.55cm, inner sep=0pt, text width={}] (k1) at (0.15,-6.05) {};
\node[key, right=0.1cm of k1] (t1) {taken unchanged};
\node[ext, minimum height=0.3cm, minimum width=0.55cm, inner sep=0pt, right=0.35cm of t1] (k2) {};
\node[key, right=0.1cm of k2] (t2) {extended and fixed};
\node[new, minimum height=0.3cm, minimum width=0.55cm, inner sep=0pt, right=0.35cm of t2] (k3) {};
\node[key, right=0.1cm of k3] (t3) {new};
\node[io, minimum height=0.3cm, minimum width=0.55cm, inner sep=0pt, right=0.35cm of t3] (k4) {};
\node[key, right=0.1cm of k4] (t4) {input or output};
\end{tikzpicture}
\end{document}
